\part{集合、简易逻辑与不等式初步}
\chapter{数理逻辑}
\section{集合}
\subsection{集合的定义、特性与表示}
在数学学习中，我们会研究一些具有共同特征的对象，我们往往称之为元素.
比如所有的质数，过某个定点的直线束......此时我们会把这些元素放在一起，得到一个集合.
集合及相关理可认为是现代数学的根基，是拥有一套公理的严格数学语言.
但是我们在高中并不会接触晦涩的公理，而是着重研究一些实用，直观的对象.
\begin{definition}\label{ref:集合}\index{集合}
    集合(set)是一些不同的、能够确定的研究对象汇集在一起构成的数学结构，其中研究的对象称之为元素(element)。
\end{definition}

我们常用大写字母$A,B,C...$表示集合，一些特殊的集合可能会使用特殊字体如$\mathcal{B},\mathbb{R},\mathfrak{F} $表示。
相对地，我们会使用小写字母$a,b,c...$来表示集合中的元素.

若元素$a$在集合$A$中，则记作\[a \in A\]

相对地，若元素$a$不在集合$A$中，则记作\[a \notin A\]

当然，如果有这样的一个集合，没有任何元素属于这个集合，那么我们称之为空集(empty set/null set)，符号为$\varnothing $.

对于两个集合$A,B$，若构成两个集合的元素相同，则记作$A,B$两个集合相等，即\[A=B\]
\begin{example}
    判断下面叙述的对象能否构成集合。若能构成集合，试图列举一些集合中的元素.
    
    $(1)$教室中比较高的同学，教室中最高的同学.

    $(2)$西游记中所有的妖怪.

    $(3)$所有的奇数，所有的偶数，所有的素数.

    $(4)$在平面上，到$A$$B$两点距离相同的所有的点.

    $(5)$单词success中的所有字母.

    $(6)$不能表示为整数或整数之比的有理数.

    $(7)$平面上所有的平行四边形.

    $(8)$除以$2$余$1$的所有整数.
\end{example}

\begin{definition}[子集]
    \[A\subseteq  B\]

    \[A \subsetneqq B\]
\end{definition}
\begin{theorem}
    \[A=B \Leftrightarrow A\subseteq B \text{且} B\subseteq A \]
\end{theorem}



知识提纲
\\集合元素的三个特性：
\\集合与元素的关系，集合与集合的关系
\\子集，空集，真子集，非空真子子集
\\常见的集合：$\mathbb{N},\mathbb{Z},\mathbb{Q},\mathbb{R},\mathbb{C}$等
\\含有n个元素的集合共有多少个子集？多少真子集？多少非空真子集？
\\集合间的运算与交并补，使用Venn图表示

\begin{example}
    \[\left\{x|y=x^2+2x+2,x\in \mathbb{R} \right\}\]
\[\left\{y|y=x^2+2x+2,x\in \mathbb{R} \right\}\]
\[\left\{(x,y)|y=x^2+2x+2,x\in \mathbb{R} \right\}\]
\[\left\{y=x^2+2x+2 \right\}\]
\end{example}

提示：字母究竟是什么意思？

\begin{example}
    \[
    \emptyset \in \left\{\emptyset\right\},
    \emptyset \subseteq \left\{\emptyset\right\},
    \emptyset \notin \emptyset,
    \emptyset \subseteq \emptyset
\]
\end{example}

\begin{example}集合A为$\left\{1,2,3,4,5,6\right\}$，集合为$\left\{4,5,6,7,8\right\}$。
则满足$S\subseteq A$且$S\cap B\neq \emptyset$的S有几个？
    
\end{example}
解：提示：将元素分为123,456,78三类。


特殊的集合表示法：区间
\\用法细则：
\\不等式的解集写集合，形如$\left\{x|-3<x<2,x \in \mathbb{R} \right\}$
\\单调区间写区间
\\取值范围尽量写区间，不要写一个光秃秃的不等式
\begin{example}
    已知 $ a \in \mathbb{R} $, 关于  $x$ 的方程 $ x^{2}+a=x  $的解组成的集合为 $ A(A \neq   \varnothing)$
    $,\left(x^{2}+a\right)^{2}+a=x $ 的解组成的集合为 $ B .$

$(1) $求证  $: A \subseteq B ;$

$(2)$ 若$  A=B$ , 求实数 $ a $ 的取值范围.
\end{example}


\clearpage
\subsection{集合的运算}
\clearpage
\hw
\begin{homework}
    \ 

    定义$1$：称一个以集合作为元素的集合为族$(collection)$.\index{族}

    定义$2$：称非空集合$X$的一个子集族$\Gamma$为X上的一个拓扑$(topology)$\index{拓扑}是指：\\
\ding{172}$ \emptyset,X\in \Gamma$;\\
\ding{173}$\Gamma$的任意多元素的并集仍在$\Gamma$中;\\
\ding{174}$\Gamma$的任意有限多元素的并集仍在$\Gamma$中.

    $(1)$族$P=\{\emptyset,X\}$，族$Q=\{A|A\subseteq X\}$,判断$P$、$Q$是否是$X$上的一个拓扑；
    
    $(2)$$\Gamma$为有限集$X$上的一个拓扑，证明：由\(\Gamma\)的所有元素关于$X$的补集构成的族$\Gamma^c$也为$X$上的一个拓扑;

    $(3)$证明：$X$为无限集合，$\tau_f=\{\complement_X A|A\text{为}X\text{的有限子集}\}\bigcup\{\varnothing\}$，则$\tau_f$为$X$上的一个拓扑.
\end{homework}



\clearpage


